\documentclass[degree-en]{wbsmba}

\title{Designing a Practical LaTeX Class for Waseda MBA Theses}
\subtitle{A teacher-acceptable approximation of the official Word template}
\author{Hikaru Waseda}
\studentid{57XXXXXX}
\zemi{Seminar on Strategic Management}
\completiondate{AY 2025}
\chiefexaminer{Prof. Kaoru Okuma}
\deputyexaminer{Prof. Tsubasa Tsubouchi}
\seconddeputyexaminer{Prof. Haruka Takata}

\begin{document}
\maketitle

\begin{abstractpage}
This sample demonstrates a practical LaTeX workflow for an MBA degree thesis at Waseda Business School.
Instead of reproducing the Word template exactly, it aims at a clean and review-friendly layout that should
be acceptable to supervisors and examiners.

The class provides a cover sheet, summary page, table of contents, main chapters, references, and appendix.
The main objective is reliable authoring rather than pixel-level fidelity.
\end{abstractpage}

\wbsfrontmatter
\tableofcontents
\wbsmainmatter

\chapter{Introduction}
\section{Background}
The official template is distributed in Word format, but a maintainable LaTeX workflow is more suitable for
long-form academic writing with repeated revisions.
That preference is consistent with established guidance on document preparation and typography
\cite{bringhurst2013elements,lamport1994latex}.

\section{Purpose}
This class offers a workable compromise between institutional appearance and standard LaTeX conventions.

\chapter{Class Design}
\section{Front Matter}
The cover sheet collects thesis metadata in a compact format inspired by the official template.

\section{Main Text}
Generous margins and double spacing are used to support annotation and review.

\chapter{Conclusion}
The result is a usable starting point for actual MBA thesis writing in the CLI environment.

\bibliographystyle{wbsplain}
\bibliography{sample-degree-en}

\appendix
\chapter{Supplementary Material}
Supplementary figures and tables may be placed here.

\end{document}
